\newpage
\section{Planos y esquemas}
\label{cap:planos_esquemas}

La correcta implementación física y lógica de una infraestructura domótica basada en inteligencia ambiental exige una rigurosa planificación espacial, eléctrica y de red. En este capítulo se presentan los esquemas gráficos y planos técnicos que definen la arquitectura de la solución propuesta. En primer lugar, se detalla el plano en planta con la distribución física de los nodos de sensado y sus áreas activas de cobertura. En segundo lugar, se exponen los esquemas electrónicos detallados para el conexionado de cada transductor con los puertos de entrada/salida (GPIO) del microcontrolador ESP32. Por último, se describe la topología y arquitectura de la red de comunicaciones locales y remotas del sistema.

\subsection{Plano de distribución de sensores en la vivienda}
\label{sec:plano_distribucion}

El sistema ha sido desplegado y evaluado en una habitación de estudio con una geometría rectangular de dimensiones aproximadas de $4.0\,\text{m} \times 3.5\,\text{m}$. La ubicación física de cada uno de los tres nodos empotrados responde a criterios de optimización de cobertura de sensado, proximidad a los elementos del entorno y minimización del ruido ambiental:

\begin{itemize}
    \item \textbf{Nodo 1 (Nodo Ambiente - $N_1$):} ubicado a la derecha del escritorio sobre la pared norte de la estancia. Esta posición elevada permite realizar lecturas térmicas, acústicas y lumínicas representativas del volumen general de la estancia sin obstrucciones.
    \item \textbf{Nodo 2 (Nodo Escritorio - $N_2$):} situado en el extremo derecho del escritorio de trabajo. Este nodo alberga el radar de ondas milimétricas (FMCW) HLK-LD2410C, orientado diagonalmente hacia la silla de estudio. Su cono de detección angular (de aproximadamente $120^\circ$) cubre la zona espacial de permanencia del usuario, permitiendo discriminar micromovimientos asociados a la respiración.
    \item \textbf{Nodo 3 (Nodo Puerta - $N_3$):} instalado directamente en la esquina inferior derecha, sobre la pared lateral derecha en el marco superior de la puerta de acceso. Gobierna el sensor de contacto de la puerta y el diodo láser emisor de la barrera óptica.
\end{itemize}

La distribución espacial de estos componentes, el mobiliario y la zona activa de influencia del transductor radar se representan vectorialmente en el plano de la estancia de la Figura~\ref{fig:plano_distribucion_planta}.

\begin{figure}[H]
    \centering
    \begin{tikzpicture}[
        wall/.style={line width=2pt, draw=black!80},
        furniture/.style={draw=black!60, fill=gray!10, thick},
        node_icon/.style={circle, draw=white, fill=Azul, text=white, font=\bfseries\small, inner sep=2.5pt},
        radar_cone/.style={draw=Verde, fill=Verde!10, fill opacity=0.3, dashed}
    ]
        % Muros exteriores (escala aprox 1.5 cm = 1 m)
        % Muro inferior continuo
        \draw[wall] (0,0) -- (6,0);
        % Muros izquierdo y superior continuos
        \draw[wall] (0,0) -- (0,5.25) -- (6,5.25);
        
        % Muro derecho con ventana en el medio y puerta en la zona inferior
        \draw[wall] (6,0) -- (6,0.2);
        % Puerta física abierta a 90 grados en la pared derecha (de y=0.2 a y=1.4)
        % Hinge at (6, 0.2), opens against the bottom wall (y=0)
        \draw[thick, draw=black!70] (6,0.2) -- (4.8,0.2);
        \draw[thin, dashed, draw=black!50] (6,1.4) arc (90:180:1.2);
        
        \draw[wall] (6,1.4) -- (6,2.2);
        % Ventana (pared derecha, de y=2.2 a y=3.5)
        \draw[thick, double, draw=black!70] (6,2.2) -- (6,3.5);
        \draw[wall] (6,3.5) -- (6,5.25);
        
        % Mobiliario: Cama pegada en la zona superior izquierda en VERTICAL
        \node[furniture, minimum width=1.35cm, minimum height=2.85cm, anchor=north west] (cama) at (0.15, 5.1) {Cama};
        % Almohada en la zona superior de la cama
        \draw[draw=black!50, fill=white] (0.25, 4.3) rectangle (1.4, 4.9);
        
        % Mesa de estudio / Escritorio a la derecha de la cama, pegado a la pared superior
        \node[furniture, minimum width=2.25cm, minimum height=1.10cm, anchor=north west] (mesa) at (1.8, 5.1) {Escritorio};
        
        % Silla de estudio debajo del escritorio
        \draw[thick, draw=black!60, fill=gray!20] (2.9, 3.2) circle (0.35);
        \node[font=\tiny] at (2.9, 3.2) {Silla};
        
        % Ubicación de los nodos (con identificador interno, sin etiquetas externas para evitar solapamientos)
        % Nodo 2 (Escritorio) en el extremo derecho del escritorio
        \node[node_icon, fill=Verde] (N2) at (3.9, 4.8) {$N_2$};
        
        % Nodo 1 (Ambiente) a la derecha del escritorio en la pared
        \node[node_icon, fill=Azul] (N1) at (5.0, 4.8) {$N_1$};
        
        % Nodo 3 (Puerta) en el marco de la puerta (pared derecha inferior)
        \node[node_icon, fill=Rojo] (N3) at (5.95, 1.4) {$N_3$};
        
        % Haz láser de barrera óptica cruzando la puerta de marco a marco
        \draw[red, dashed, thick] (5.95, 1.4) -- (5.95, 0.2);
        \node[red, font=\tiny, anchor=west] at (5.95, 0.8) {Láser};
        
        % Cono del radar mmWave (120º desde el nodo N2 diagonalmente hacia la Silla)
        \begin{scope}
            \clip (0,0) rectangle (6,5.25);
            \draw[radar_cone] (3.9, 4.8) -- +(175:2.8) arc (175:295:2.8) -- cycle;
        \end{scope}
        \node[font=\tiny\bfseries, color=Verde!80!black, align=center] at (2.0, 2.5) {Cobertura Radar\\mmWave (120º)};
        
        % Leyenda de los Nodos fuera del plano (horizontal, debajo de la estancia - ampliado a 10cm para evitar solapamientos)
        \draw[draw=black!30, fill=white, rounded corners=2pt, thick] (-2.0, -0.8) rectangle (8.0, -0.2);
        
        \node[circle, draw=white, fill=Azul, text=white, font=\bfseries\tiny, inner sep=1.5pt] (leg_n1) at (-1.5, -0.5) {$N_1$};
        \node[anchor=west, font=\tiny] at (-1.2, -0.5) {Nodo Ambiente};
        
        \node[circle, draw=white, fill=Verde, text=white, font=\bfseries\tiny, inner sep=1.5pt] (leg_n2) at (2.0, -0.5) {$N_2$};
        \node[anchor=west, font=\tiny] at (2.3, -0.5) {Nodo Escritorio};
        
        \node[circle, draw=white, fill=Rojo, text=white, font=\bfseries\tiny, inner sep=1.5pt] (leg_n3) at (5.5, -0.5) {$N_3$};
        \node[anchor=west, font=\tiny] at (5.8, -0.5) {Nodo Puerta};
    \end{tikzpicture}
    \caption{Plano en planta de la habitación de estudio, mobiliario y distribución física de los nodos de control.}
    \label{fig:plano_distribucion_planta}
\end{figure}

\subsection{Esquemas electrónicos de los nodos}
\label{sec:esquemas_electronicos}

La interconexión eléctrica de los componentes electrónicos y transductores con las placas de desarrollo ESP32 se ha estructurado de modo que se respeten estrictamente las especificaciones de tensión e inmunidad al ruido. Cada nodo de hardware opera con una tensión lógica máxima de $3.3\,\text{V}$ en sus puertos GPIO, implementando circuitos acondicionadores cuando es preciso.

En la Figura~\ref{fig:esquema_amb_puerta} se detallan las conexiones electrónicas correspondientes al Nodo 1 (Ambiente) y al Nodo 3 (Puerta). El esquema se ha organizado en dos secciones verticales de gran amplitud, asignando una separación de $4.5\,\text{cm}$ en el eje horizontal y de $1.7\,\text{cm}$ en el eje vertical entre sensores, lo que elimina cualquier solapamiento en las etiquetas de los pines.

\begin{figure}[H]
    \centering
    \begin{tikzpicture}[
        block/.style={rectangle, draw=black!80, fill=gray!5, thick, minimum width=2.4cm, minimum height=3.0cm, align=center, font=\bfseries\small},
        sensor/.style={rectangle, draw=Azul!80, fill=Azul!5, thick, minimum width=2.1cm, minimum height=0.9cm, align=center, font=\footnotesize},
        actuator/.style={rectangle, draw=Rojo!80, fill=Rojo!5, thick, minimum width=2.1cm, minimum height=0.9cm, align=center, font=\footnotesize},
        arrow/.style={thick, ->, >=stealth, draw=black!70}
    ]
        % ================= NODO 1: AMBIENTE (PARTE SUPERIOR) =================
        \node[block, draw=Azul!90, fill=Azul!2] (esp_amb) at (0, 2.2) {ESP32\\Nodo 1\\(Ambiente)};
        
        \node[sensor] (dht) at (-5.5, 4.2) {DHT11\\(Clima)};
        \node[sensor] (ldr) at (-5.5, 2.2) {Circuito LDR\\(Analógico)};
        \node[sensor] (sound) at (-5.5, 0.2) {Sensor Sonido\\(Big Sound)};
        \node[actuator] (ssr) at (5.5, 2.2) {Relé SSR\\(Actuador)};
        
        \draw[arrow] (dht.east) -- node[above=3pt, midway, font=\tiny\bfseries] {GPIO 4} (esp_amb.150);
        \draw[arrow] (ldr.east) -- node[above=3pt, midway, font=\tiny\bfseries] {GPIO 32} (esp_amb.180);
        \draw[arrow] (sound.east) -- node[above=3pt, midway, font=\tiny\bfseries] {GPIO 27} (esp_amb.210);
        \draw[arrow] (esp_amb.0) -- node[above=3pt, midway, font=\tiny\bfseries] {GPIO 26} (ssr.west);
        
        \node[anchor=west, font=\bfseries\small, color=Azul!80!black] at (-5.5, 5.4) {NODO 1: MONITORIZACIÓN AMBIENTAL};
        
        % Línea de separación estética
        \draw[dashed, gray!45, thick] (-6.0, -0.6) -- (6.0, -0.6);
        
        % ================= NODO 3: PUERTA (PARTE INFERIOR) =================
        \begin{scope}[yshift=-3.8cm]
            \node[block, draw=Rojo!90, fill=Rojo!2] (esp_puerta) at (0, 0) {ESP32\\Nodo 3\\(Puerta)};
            
            \node[sensor] (reed) at (-5.5, 0) {Interruptor Reed\\(Contacto)};
            
            \draw[arrow] (reed.east) -- node[above=3pt, midway, font=\tiny\bfseries] {GPIO 4 (Pull-up)} (esp_puerta.180);
            
            \node[anchor=west, font=\bfseries\small, color=Rojo!80!black] at (-5.5, 2.2) {NODO 3: SEGURIDAD DE ACCESO};
        \end{scope}
    \end{tikzpicture}
    \caption{Esquema de conexionado electrónico y puertos GPIO para el Nodo 1 (Ambiente) y el Nodo 3 (Puerta) con distribución vertical espaciada anti-solapamiento.}
    \label{fig:esquema_amb_puerta}
\end{figure}

Por otro lado, la comunicación del Nodo 2 (Escritorio) con el radar mmWave HLK-LD2410C se implementa mediante comunicación serie UART. En la Figura~\ref{fig:esquema_uart_escritorio} se detalla este circuito. Se ha incrementado la separación horizontal entre los bloques a $7.5\,\text{cm}$, ofreciendo un espacio de holgura de $4.7\,\text{cm}$ en el centro del bus serie que permite representar las etiquetas de transmisión (TX2) y recepción (RX2) de forma clara y sin colisiones con los bordes de los dispositivos.

\begin{figure}[H]
    \centering
    \begin{tikzpicture}[
        block/.style={rectangle, draw=black!80, fill=gray!5, thick, minimum width=2.8cm, minimum height=3.2cm, align=center, font=\bfseries\small},
        radar/.style={rectangle, draw=Verde!80, fill=Verde!5, thick, minimum width=2.8cm, minimum height=3.2cm, align=center, font=\bfseries\small},
        connect/.style={thick, draw=black!80}
    ]
        % Bloques con separación horizontal aumentada a 7.5cm
        \node[block, draw=Azul!90, fill=Azul!2] (esp) at (0,0) {ESP32\\Nodo 2\\(Escritorio)};
        \node[radar] (rad) at (7.5,0) {Radar mmWave\\HLK-LD2410C};
        
        % Conexión 5V (Ruta superior externa)
        \draw[connect, color=red, thick] (esp.north) -- ++(0,0.9) -| node[above, pos=0.5, font=\footnotesize\bfseries] {5V / VIN (Alimentación)} (rad.north);
        
        % Conexión GND (Ruta inferior externa, etiqueta por encima de la línea)
        \draw[connect, color=black, thick] (esp.south) -- ++(0,-0.9) -| node[above, pos=0.5, font=\footnotesize\bfseries] {GND (Tierra Común)} (rad.south);
        
        % Buses de comunicación serie cruzados en la zona central con espacio ampliado
        \draw[connect, ->, >=stealth, color=blue, thick] (esp.15) -- node[above, midway, font=\footnotesize\bfseries, fill=white, inner sep=2pt] {TX2 (GPIO 17) $\rightarrow$ RX} (rad.165);
        \draw[connect, <-, >=stealth, color=orange!80!black, thick] (esp.-15) -- node[below, midway, font=\footnotesize\bfseries, fill=white, inner sep=2pt] {RX2 (GPIO 16) $\leftarrow$ TX} (rad.195);
        
        % Caja de anotación (desplazada por debajo de la línea GND para evitar cruces)
        \node[anchor=north, align=center, font=\scriptsize\itshape, draw=gray!30, fill=gray!5, rounded corners=2pt, inner sep=5pt] at (3.75,-2.9) {
            Tasa de transferencia: 256\,000 baudios\\
            Nivel de tensión de señal: 3.3\,V TTL (Compatibilidad directa)
        };
    \end{tikzpicture}
    \caption{Esquema detallado de la conexión serie UART y alimentación entre el ESP32 y el radar HLK-LD2410C con espacio de bus ampliado.}
    \label{fig:esquema_uart_escritorio}
\end{figure}

\subsection{Topología de la red de comunicaciones}
\label{sec:topologia_red}

El intercambio de información y la orquestación contextual se estructuran sobre una topología en estrella segmentada localmente que se conecta de forma segura con el exterior mediante una red privada virtual (VPN). Los flujos de comunicación se dividen en tres planos lógicos diferenciados:

\begin{enumerate}
    \item \textbf{Red Troncal Local (WLAN Wi-Fi):} los tres nodos de control basados en ESP32 se asocian al punto de acceso local (router doméstico) mediante comunicación inalámbrica en la banda de $2.4\,\text{GHz}$. Para agilizar el proceso de descubrimiento e interconexión y mitigar la latencia, cada microcontrolador tiene configurada una dirección IP estática dentro del rango de subred local (\texttt{192.168.1.X}).
    \item \textbf{Canal de Orquestación perimetral (ESPHome API):} el servidor central (Raspberry Pi 4) está conectado físicamente al router mediante un enlace Ethernet Gigabit Cat6 para asegurar una conexión robusta y libre de colisiones. El servidor aloja los contenedores aislados mediante Docker. Home Assistant establece conexiones de sockets TCP directas (puerto 6053) con la API nativa de cada microcontrolador para la transmisión periódica de telemetría y la recepción de comandos de conmutación.
    \item \textbf{Canal de Telemetría Remota Segura (Tailscale VPN):} con el fin de habilitar el acceso remoto y el control del panel Lovelace desde el exterior de la vivienda sin vulnerar la red local (evitando la redirección de puertos entrantes en el router), se despliega un túnel de cifrado extremo a extremo basado en la VPN Tailscale (protocolo WireGuard). Tanto el servidor Raspberry Pi como los terminales autorizados de administración se integran en la misma malla segura de red privada (\textit{Tailnet}), garantizando el acceso cifrado mediante protocolo HTTPS.
\end{enumerate}

La arquitectura lógica de red del sistema se ilustra en la Figura~\ref{fig:topologia_red_sistema}. Para evitar solapamientos y cruces de líneas sobre los bloques de red, se ha diseñado una topología de flujo ordenada izquierda-derecha con enrutamientos perimetrales superior e inferior. El router de acceso Wi-Fi se ubica a la izquierda, la columna de microcontroladores se sitúa en el centro, y los elementos de procesamiento en el borde y acceso VPN ocupan el flanco derecho del diagrama.

\begin{figure}[H]
    \centering
    \shorthandoff{<}\shorthandoff{>}
    \resizebox{\textwidth}{!}{%
    \begin{tikzpicture}[
        device/.style={rectangle, draw=black!80, fill=gray!5, thick, minimum width=2.2cm, minimum height=1.0cm, align=center, font=\bfseries\scriptsize},
        srv/.style={rectangle, draw=Azul!90, fill=Azul!5, thick, minimum width=3.5cm, minimum height=2.6cm, align=center, font=\bfseries\scriptsize},
        rtr/.style={circle, draw=Naranja!80, fill=Naranja!10, thick, minimum size=1.8cm, align=center, font=\bfseries\scriptsize},
        cloud_node/.style={rectangle, rounded corners=6pt, draw=gray!50, fill=gray!5, thick, minimum width=2.2cm, minimum height=1.2cm, align=center, font=\bfseries\scriptsize},
        eth/.style={thick, draw=blue!80, line width=1.5pt},
        wifi/.style={thick, draw=green!60!black, dashed, line width=1.2pt},
        tunnel/.style={draw=red!70, line width=1.8pt, dashed, arrows={stealth-stealth}, shorten >=2pt, shorten <=2pt},
        line/.style={thick, draw=black!60}
    ]%
        % ================= NODO DE APORTACIÓN FÍSICA E INSTALACIONES (EJE CENTRAL) =================
        % 1. Enrutador local (a la izquierda de la escena)
        \node[rtr] (router) at (-6.5, 0) {Router\\Doméstico\\(192.168.1.1)};
        % 2. Tres nodos sensores en la columna central
        \node[device] (n1) at (-1.5, 1.8) {Nodo 1 (Ambiente)\\ESP32\\IP: 192.168.1.150};
        \node[device] (n2) at (-1.5, 0) {Nodo 2 (Escritorio)\\ESP32\\IP: 192.168.1.151};
        \node[device] (n3) at (-1.5, -1.8) {Nodo 3 (Puerta)\\ESP32\\IP: 192.168.1.152};
        % 3. Servidor central en el borde (a la derecha)
        \node[srv] (rpi) at (4.5, -0.6) {Raspberry Pi 4 (Servidor)\\IP Local: 192.168.1.100\\[0.2cm]
            \parbox{2.9cm}{\centering\color{black!80}\bfseries\tiny
            \begin{tabular}{|c|}
            \hline
            \textbf{Entorno Docker} \\
            Home Assistant (8123) \\
            ESPHome API (6053) \\
            SQLite DB \\
            Tailscale VPN \\
            \hline
            \end{tabular}%
            }%
        };
        % ================= CONEXIONES WI-FI (LADO IZQUIERDO Y DIAGONALES) =================
        \draw[wifi] (router) -- node[above left, font=\tiny\bfseries, color=green!50!black] {Wi-Fi} (n1);
        \draw[wifi] (router) -- node[above, font=\tiny\bfseries, color=green!50!black] {Wi-Fi} (n2);
        \draw[wifi] (router) -- node[below left, font=\tiny\bfseries, color=green!50!black] {Wi-Fi} (n3);
        % ================= CONEXIONES DE API LÓGICA (LADO DERECHO) =================
        \draw[thick, arrows={stealth-stealth}, color=purple!80, line width=1.0pt] (n1.east) to[bend left=5] node[above, midway, sloped, font=\tiny\bfseries] {API (TCP 6053)} (rpi.130);
        \draw[thick, arrows={stealth-stealth}, color=purple!80, line width=1.0pt] (n2.east) -- (rpi.170);
        \draw[thick, arrows={stealth-stealth}, color=purple!80, line width=1.0pt] (n3.east) to[bend right=5] node[below, midway, sloped, font=\tiny\bfseries] {API (TCP 6053)} (rpi.210);
        % ================= CONEXIÓN CABLEADA E INTERNET (BLES Y RUTA PERIMETRAL) =================
        % Enlace Ethernet Gigabit (Ruta inferior perimetral por debajo de n3, pos=0.35 para centrar en el espacio)
        \draw[eth] (router.south) -- ++(0, -1.6) -| node[above, pos=0.35, font=\tiny\bfseries, color=blue!80!black] {Ethernet Cat6} (rpi.south);
        % Conexión lógica WAN a Internet (Ruta superior perimetral por encima de n1, pos=0.35 para centrar en el espacio)
        \node[cloud_node] (cloud) at (4.5, 1.8) {Internet\\(Red Pública)};
        \draw[line] (router.north) -- ++(0, 1.8) -| node[above, pos=0.35, font=\tiny\bfseries] {Internet WAN} (cloud.north);
        % Conexión directa Servidor - Nube (línea vertical sin cruces)
        \draw[line] (rpi.north) -- (cloud.south);
        % Terminal de Administración de Cliente (más a la derecha para dejar espacio)
        \node[device] (client) at (9.5, 0.6) {Dispositivo Cliente\\(Smartphone / PC)\\IP VPN: 100.75.10.20};
        % Enlace del Cliente a Internet
        \draw[line] (client.north) |- (cloud.east);
        % Túnel seguro VPN Tailscale (Ruta por flanco derecho, etiqueta centrada debajo de la sección horizontal)
        \draw[tunnel] (client.south) -- ++(0, -0.7) -- node[below, midway, font=\tiny\bfseries, color=red!80!black] {VPN Tailscale} (rpi);
    \end{tikzpicture}%
    }
    \shorthandon{<}\shorthandon{>}
    \caption{Topología general de red del sistema domótico perimetral y túnel lógico de acceso remoto con desacoplamiento espacial de flujos y control anti-cruce.}
    \label{fig:topologia_red_sistema}
\end{figure}